\begin{figure}[t]
\centering
\begin{tikzpicture}
\begin{axis}[
  width=0.92\linewidth, height=4.8cm,
  xbar, bar width=11pt,
  xmin=-14, xmax=30,
  xlabel={MaxPressure vs fixed-time (\%); positive = MaxPressure worse},
  xlabel style={font=\sffamily\scriptsize},
  symbolic y coords={Euston cal., Euston, Bloomsbury, Old Street, Bloomsbury cal.},
  ytick=data, y tick label style={font=\sffamily\scriptsize},
  xtick={-10,0,10,20,30}, xticklabel style={font=\sffamily\scriptsize},
  axis line style={nbgrey}, tick style={nbgrey},
  xmajorgrids, grid style={nbgrey!25},
  enlarge y limits=0.16,
  nodes near coords, every node near coord/.append style={font=\sffamily\tiny, /pgf/number format/fixed, /pgf/number format/precision=1},
]
\addplot[fill=nbblue!55, draw=nbblue] coordinates
  {(14.2,Euston cal.) (22.3,Euston) (0.2,Bloomsbury) (-8.2,Old Street) (-8.8,Bloomsbury cal.)};
\draw[nbgrey, dashed] (axis cs:0,Euston cal.) -- (axis cs:0,Bloomsbury cal.);
\end{axis}
\end{tikzpicture}
\caption{The classical baseline is not uniformly strong: MaxPressure is
significantly worse than naive fixed-time on the Euston corridor,
significantly better on Old Street, and tied on the oversaturated
Bloomsbury grid; calibrated to clear, the grid favours MaxPressure and the
corridor stays inconclusive. Every controller claim is therefore stated
against the fixed-time floor with MaxPressure alongside.}
\label{fig:signflip}
\end{figure}
